% =====================================================================
% preamble.tex -- shared look and helper commands for the B3 slide deck
% Owner: T3 (slides lead). Change this file only after telling the team:
% every section file depends on the colours and macros below.
% Compile with pdfLaTeX (Overleaf default).
% =====================================================================

\usepackage[T1]{fontenc}
\usepackage{lmodern}
\usepackage{amsmath,amssymb}
\usepackage{booktabs}
\usepackage{array}
\usepackage{tikz}
\usetikzlibrary{arrows.meta,positioning,calc,shapes.geometric,fit,backgrounds}
\usepackage{pgfplots}
\pgfplotsset{compat=1.18}
\usepackage{listings}
\usepackage[most]{tcolorbox}
\usepackage{appendixnumberbeamer}

% ---------------------------------------------------------------------
% Colour palette (light theme)
% ---------------------------------------------------------------------
\definecolor{kdnavy}{HTML}{1F3A5F}    % headings
\definecolor{kdblue}{HTML}{2F6DB5}    % structure, main data series
\definecolor{kdpale}{HTML}{EAF1FB}    % title band, block titles
\definecolor{kdpaler}{HTML}{F5F8FD}   % block bodies
\definecolor{kdteal}{HTML}{178F87}    % second data series, "Ours"
\definecolor{kdorange}{HTML}{D9731A}  % accent, baselines, to-do boxes
\definecolor{kdred}{HTML}{C0392B}     % errors, warnings
\definecolor{kdgray}{HTML}{5B6575}    % muted text
\definecolor{kdlight}{HTML}{F2F4F7}   % footline
\definecolor{kdink}{HTML}{1E2430}     % body text
\definecolor{kdcode}{HTML}{F6F8FA}    % code background
\definecolor{kdtodo}{HTML}{FFF6E8}    % to-do box background

% ---------------------------------------------------------------------
% Beamer colours, fonts and templates
% ---------------------------------------------------------------------
\setbeamertemplate{navigation symbols}{}
\setbeamercolor{normal text}{fg=kdink,bg=white}
\setbeamercolor{structure}{fg=kdblue}
\setbeamercolor{alerted text}{fg=kdorange}
\setbeamercolor{frametitle}{fg=kdnavy,bg=kdpale}
\setbeamercolor{framesubtitle}{fg=kdgray}
\setbeamercolor{title}{fg=kdnavy}
\setbeamercolor{subtitle}{fg=kdblue}
\setbeamercolor{block title}{fg=kdnavy,bg=kdpale}
\setbeamercolor{block body}{fg=kdink,bg=kdpaler}
\setbeamercolor{block title alerted}{fg=white,bg=kdorange}
\setbeamercolor{block body alerted}{fg=kdink,bg=kdtodo}
\setbeamercolor{block title example}{fg=white,bg=kdteal}
\setbeamercolor{block body example}{fg=kdink,bg=kdteal!6}
\setbeamercolor{footline}{fg=kdgray,bg=kdlight}
\setbeamercolor{item}{fg=kdblue}
\setbeamercolor{section in toc}{fg=kdnavy}
\setbeamercolor{section in toc shaded}{fg=kdgray}

\setbeamerfont{frametitle}{size=\large,series=\bfseries}
\setbeamerfont{framesubtitle}{size=\footnotesize,series=\mdseries}
\setbeamerfont{footline}{size=\tiny}
\setbeamerfont{block title}{size=\small,series=\bfseries}

\setbeamertemplate{itemize items}[circle]
\setbeamertemplate{enumerate items}[default]
\setbeamertemplate{blocks}[rounded]
\setbeamertemplate{section in toc}[sections numbered]
\setbeamertemplate{section in toc shaded}[default][40]

% Frame title: a pale band across the slide with a thin blue rule below.
\makeatletter
\setbeamertemplate{frametitle}{%
  \nointerlineskip
  \begin{beamercolorbox}[wd=\paperwidth,leftskip=0.5cm,rightskip=0.5cm,sep=0.2cm]{frametitle}%
    \usebeamerfont{frametitle}\strut\insertframetitle\strut
    \ifx\insertframesubtitle\@empty\else
      \par{\usebeamerfont{framesubtitle}\usebeamercolor[fg]{framesubtitle}\insertframesubtitle\strut\par}%
    \fi
  \end{beamercolorbox}%
  \nointerlineskip
  \hbox{\color{kdblue}\rule{\paperwidth}{0.6pt}}%
  \vskip0.1cm
}
\makeatother

% Footline: project name, current section, slide number.
\setbeamertemplate{footline}{%
  \leavevmode%
  \hbox{%
    \begin{beamercolorbox}[wd=.40\paperwidth,ht=2.4ex,dp=1ex,leftskip=0.5cm]{footline}%
      \usebeamerfont{footline}Project B3 $\cdot$ Knowledge Distillation
    \end{beamercolorbox}%
    \begin{beamercolorbox}[wd=.40\paperwidth,ht=2.4ex,dp=1ex,center]{footline}%
      \usebeamerfont{footline}\insertsectionhead
    \end{beamercolorbox}%
    \begin{beamercolorbox}[wd=.20\paperwidth,ht=2.4ex,dp=1ex,right,rightskip=0.5cm]{footline}%
      \usebeamerfont{footline}\insertframenumber{} / \inserttotalframenumber
    \end{beamercolorbox}}%
  \vskip0pt%
}

% Section divider: big number, section title, presenter and minutes,
% and the roadmap with the current section highlighted.
% Each section file sets \sectioninfo just before its \section command.
\newcommand{\sectioninfo}{}
\AtBeginSection[]{%
  \begin{frame}[plain,noframenumbering]
    \begin{tikzpicture}[remember picture,overlay]
      \fill[kdpaler] (current page.north west) rectangle (current page.south east);
      \fill[kdblue] (current page.north west) rectangle ([xshift=0.22cm]current page.south west);
    \end{tikzpicture}%
    \begin{columns}[c]
      \begin{column}{0.52\textwidth}
        {\color{kdblue}\fontsize{38}{40}\selectfont\bfseries\thesection\par}
        \vspace{0.2cm}
        {\color{kdnavy}\Large\bfseries\insertsectionhead\par}
        \vspace{0.3cm}
        {\color{kdgray}\normalsize\sectioninfo\par}
      \end{column}
      \begin{column}{0.42\textwidth}
        \scriptsize\tableofcontents[currentsection,hideallsubsections]
      \end{column}
    \end{columns}
  \end{frame}
}

% ---------------------------------------------------------------------
% Helper commands used on the slides
% ---------------------------------------------------------------------
% Small coloured labels for frame titles:
%   \est  = claim taken from the literature
%   \ours = result from our own runs (a file in results/)
%   \who{T3} = who presents this slide
\newcommand{\badge}[2]{{\setlength{\fboxsep}{1.6pt}\colorbox{#1}{\textcolor{white}{\scriptsize\bfseries\strut #2}}}}
\newcommand{\est}{\badge{kdblue}{Established}}
\newcommand{\ours}{\badge{kdteal}{Ours}}
\newcommand{\who}[1]{\badge{kdgray}{#1}}
% A source in grey after a statement: \src{Hinton et al., 2015}
\newcommand{\src}[1]{{\color{kdgray}\scriptsize(#1)}}
% Arm names in typewriter type
\newcommand{\arm}[1]{\texttt{#1}}
% A path or file in the repository (not \path: TikZ already uses that name)
\newcommand{\repofile}[1]{{\color{kdnavy}\texttt{#1}}}

% To-do box: marks content that is missing from the repository or still
% has to be written. Replace the box with real content when it exists.
%   \begin{todobox}{Seed 1 results (T1)} what is needed ... \end{todobox}
\newtcolorbox{todobox}[1]{enhanced,colback=kdtodo,colframe=kdorange,
  boxrule=0.6pt,arc=2pt,left=5pt,right=5pt,top=3pt,bottom=3pt,
  coltitle=white,colbacktitle=kdorange,fonttitle=\bfseries\scriptsize,
  title={To add: #1},fontupper=\scriptsize}

% Info box for short key messages
\newtcolorbox{keybox}{enhanced,colback=kdpale,colframe=kdblue,
  boxrule=0.6pt,arc=2pt,left=6pt,right=6pt,top=4pt,bottom=4pt,fontupper=\small}

% ---------------------------------------------------------------------
% TikZ helpers
% ---------------------------------------------------------------------
% "visible on=<2->" shows a TikZ element from overlay 2 onwards while
% keeping its space reserved, so the picture does not jump.
\tikzset{
  invisible/.style={opacity=0,text opacity=0},
  visible on/.style={alt={#1{}{invisible}}},
  alt/.code args={<#1>#2#3}{\alt<#1>{\pgfkeysalso{#2}}{\pgfkeysalso{#3}}},
  stage/.style={draw=kdblue,fill=kdpale,rounded corners=3pt,align=center,
    font=\scriptsize,text width=1.85cm,minimum height=1.15cm,inner sep=3pt},
  stageT/.style={stage,draw=kdteal,fill=kdteal!8},
  stageX/.style={stage,draw=kdorange,fill=kdtodo,dashed},
  flow/.style={-{Stealth[length=2mm]},thick,draw=kdgray},
  filelabel/.style={font=\tiny\ttfamily,text=kdgray,align=center},
}

% pgfplots defaults for this deck
\pgfplotsset{
  kdaxis/.style={
    axis line style={kdgray},
    tick style={kdgray},
    tick label style={font=\scriptsize,text=kdink},
    label style={font=\scriptsize,text=kdink},
    legend style={font=\scriptsize,draw=none,fill=none},
    title style={font=\small\bfseries,text=kdnavy},
    major grid style={kdlight},
  },
}

% ---------------------------------------------------------------------
% Code listings
% ---------------------------------------------------------------------
\lstdefinestyle{kdpy}{
  language=Python,
  basicstyle=\ttfamily\scriptsize,
  keywordstyle=\color{kdblue}\bfseries,
  commentstyle=\color{kdgray}\itshape,
  stringstyle=\color{kdteal},
  backgroundcolor=\color{kdcode},
  frame=single,framerule=0.4pt,rulecolor=\color{kdpale},
  xleftmargin=3pt,xrightmargin=3pt,
  showstringspaces=false,columns=fullflexible,keepspaces=true,
  breaklines=true,tabsize=4,aboveskip=3pt,belowskip=3pt,upquote=true,
}
\lstdefinestyle{kdplain}{
  style=kdpy,language={},
  morecomment=[l]{\#},
}
\lstset{style=kdpy}
